\section{Fluid Model and Equilibrium}\label{sec:fluid}
Section~\ref{sec:model} identifies the central admission-control trade-off: the scheduler should admit enough prompts to exploit batching, but not so many that their KV caches later overflow memory. This section derives the balanced operating point for that trade-off. The key state variable is the distribution of resident prompts across prefill and decode stages. As prompts move deeper into decode, each prompt carries a larger KV cache; if too many prompts accumulate at late stages, the system can run out of memory and restart work through eviction.

We develop a fluid model that replaces stochastic sample paths with average flows and characterizes an ideal-fluid equilibrium in which arrivals, stage transitions, and completions occur at matching rates. For a given arrival vector, this equilibrium gives two reference quantities: the memory requirement \(M^*\) needed to support the equilibrium and the effective throughput \(\throughput^*\) achieved when the offered load can be completed without eviction. Equivalently, for a fixed memory capacity \(C\), the ideal-fluid stability region consists of arrival vectors whose equilibrium memory requirement satisfies \(M^*\le C\). We derive these quantities first for a single prompt type, which isolates the stage-by-stage memory accumulation, and then for multiple prompt types sharing one memory constraint. The same equilibrium also supplies the design target for WAIT in Section~\ref{sec:wait}: keep the in-service population close to the implied allocation across decode stages, thereby avoiding eviction-induced restarts while approaching the largest throughput attainable under the offered load and memory constraint.

\subsection{Single-Type Fluid Model}
Consider a single prompt type \(j\). A resident prompt uses \(l_j\) units of KV cache immediately after prefill and \(l_j+s\) units at decode stage \(s=1,\dots,l_j'\); after stage \(l_j'\), the prompt completes and releases its cache. Let \(n_j^*\) denote the equilibrium number of type-\(j\) prompts at each stage. The balanced state therefore has \(n_j^*\) prompts in each of the \(l_j'+1\) stages.

Before deriving the equilibrium, we identify the load condition under which batching can be effective. Adding prompts to a batch is useful only if the additional completions outweigh the longer iteration time. Across its prefill and decode stages, a type-\(j\) prompt contributes memory-dependent work proportional to \((l_j' + 1)(l_j + l_j'/2)\). If arrivals are high enough that this memory-dependent work alone saturates the server, batching cannot restore stability.

\begin{proposition}\label{prop:large_rate}
When the condition $\lambda_j d_1 (l_j' + 1) \left( l_j + \frac{l_j'}{2} \right) \geq 1$ is satisfied, the system is unstable in the sense that the expected average latency under any scheduling policy with arrival rate \( \lambda_j \) grows linearly with time. Formally,
$$
\mathbb{E}[\textbf{Latency}^{(T,\pi)}] = \Omega(T) \quad \text{as} \quad T \to \infty.
$$
\end{proposition}

The proof compares arrivals and completions in a balanced batch and shows that the arrival-completion gap remains positive when the condition holds (Appendix~\ref{appendix:proof_large_rate}). We therefore focus on the nontrivial regime $\lambda_j d_1 (l_j' + 1) \left( l_j + \frac{l_j'}{2} \right) <1$, where batching can support an ideal-fluid equilibrium.

At equilibrium, the \(l_j'+1\) stages each contain \(n_j^*\) prompts, so the memory usage is
\begin{equation}
    M_j^* = n_j^* \sum_{s=0}^{l_j^{\prime}}(l_j+s) = n_j^*(l_j'+1)\prn{l_j+\frac{l_j'}{2}},
\end{equation}
where \(l_j+l_j'/2\) is the average KV-cache usage of a type-\(j\) prompt across its prefill and decode stages.

Flow balance requires arrivals during one iteration to match completions. Since one iteration takes time \(d_0+d_1M_j^*\), \(\lambda_j(d_0+d_1M_j^*)\) prompts arrive and \(n_j^*\) prompts complete, giving
\begin{equation}
\label{eq:single_type_equilibrium}
\underset{\text{Time per iteration}}{\underbrace{(d_0+d_1M_j^{*} )}}\lambda_j =\left(d_0+d_1 n_j^*(l_j'+1)\left( l_j +\frac{l_j^{\prime}}{2} \right) \right)\lambda_j= n_j^*.
\end{equation}

Solving gives \(n_j^* = d_0\lambda_j/[1-d_1\lambda_j(l_j^{\prime}+1)(l_j+l_j^{\prime}/2)]\). Substituting this value into the memory expression yields
\begin{equation}
    \label{eq:memory_single}
M_j^* = n_j^*(l_j'+1)\left( l_j+\frac{l_j^{\prime}}{2} \right)
=  \frac{d_0\lambda_j(l_j^{\prime}+1)\left( l_j+\frac{l_j^{\prime}}{2} \right) }{1-d_1\lambda_j(l_j^{\prime}+1)(l_j+\frac{l_j^{\prime}}{2})}.
\end{equation}

The denominator is the slack in the batching feasibility condition; as it approaches zero, the inventory required for balance diverges. The balanced throughput is completed decode tokens per iteration divided by iteration time:
\begin{equation}\label{eq:throughput_fluid}
\throughput_j^* = \frac{n_j^* \cdot l_j'}{d_0+d_1n_j^*(l_j'+1)\prn{l_j+\frac{l_j'}{2}}} = \lambda_j l_j',
\end{equation}
where the equality follows from the flow-balance condition \eqref{eq:single_type_equilibrium}.

\subsection{Multiple-Type Fluid Model}
We now extend the balance calculation to \(m\) prompt types. Type \(j\) has input length \(l_j\), decode length \(l_j'\), and arrival rate \(\lambda_j\). Because all types draw from the same memory pool, the equilibrium is determined by the joint allocation of active prompts across decode stages and the aggregate KV-cache usage it induces.

The balance equations are written at the scheduler's iteration scale: an iteration selects a batch, advances active prompts by one stage, and updates their KV-cache usage. Let \(n_j^*\) denote the equilibrium number of type-\(j\) prompts at each stage. The resulting memory usage is
\begin{equation}
\label{eq:memory_multi1}
    M^* = \sum_{j=1}^m n_j^*(l_j'+1)\left(l_j + \frac{l_j^\prime}{2}\right)
\end{equation}
Flow balance requires \((d_0+d_1M^*)\lambda_j = n_j^*\) for each type \(j\), because \(n_j^*\) type-\(j\) prompts complete in one iteration. Substituting these identities into \eqref{eq:memory_multi1} gives
\begin{equation}
\label{eq:memory_multi2}
M^* = \frac{d_0\sum_{j=1}^m \lambda_j (l_j^{\prime}+1)\left(l_j + \frac{l_j^{\prime}}{2} \right)}{1-d_1\sum_{j=1}^m \lambda_j (l_j^{\prime}+1)\left(l_j + \frac{l_j^{\prime}}{2} \right) } 
\end{equation}
The corresponding iteration time is
\begin{equation}    \label{eq:time_fluid}
    \begin{aligned}
            &\Delta T^* = d_0 + d_1 M^*  \\
                  &= \frac{d_0}{1 - d_1 \sum_{j=1}^m \lambda_j (l_j' + 1)(l_j + \tfrac{1} {2}l_j')},
    \end{aligned}
\end{equation}
and the fluid throughput benchmark is
\begin{equation}
    \label{throughput_fluid_multiple}
    \begin{aligned}
        &\throughput^*
        = \sum_{j=1}^m \frac{n_j^* \cdot l_j^{\prime}}{d_0 + d_1 M^*}
        = \sum_{j=1}^m \lambda_j l_j^{\prime}
    \end{aligned}
\end{equation}
The ideal fluid model gives an upper bound for any online policy:
% Detailed proofs are presented in Appendix~\ref{appendix::proof online policy expected throughput upper bound, multiple type}.

\begin{proposition}\label{prop::online policy expected throughput upper bound, multiple type}
Suppose that the memory capacity satisfies \(C \geq M^*\), where \(M^*\) denotes the memory requirement needed to support the ideal-fluid equilibrium in \eqref{eq:memory_multi2}. Then, for any online policy \(\pi \in \Pi\), the expected throughput satisfies
\[
\mathbb{E}[\throughput^{(T,\pi)}] \leq \throughput^*,
\]
where \(\throughput^*\) is given in \eqref{throughput_fluid_multiple}.
\end{proposition}
The fluid model serves two purposes. First, it provides the throughput benchmark in Proposition~\ref{prop::online policy expected throughput upper bound, multiple type}. Second, the ideal-fluid equilibrium identifies the operating point a policy should maintain: roughly \(n_j^*\) type-\(j\) prompts at each decode stage, with total memory requirement \(M^*\). Accumulating too much late-stage work consumes memory and risks eviction, while holding too little late-stage work leaves completions below the fluid rate. Section~\ref{sec:wait} uses \(M^*\) and \(n_j^*\) to define thresholds that keep the decode-stage populations near the ideal-fluid equilibrium.

% The \textit{iteration time} \(\tau^*\), the duration to process one batch, is:
% \[
% \tau^* = \max \left\{ d_f \cdot \max_j \{ l_j \}, \; d_1 \cdot \sum_{j=1}^m n_{j, \text{decode}}^* \left( l_j + \frac{l_j' + 1}{2} \right) \right\},
% \]
% where \(d_f\) and \(d_1\) are prefill and decode processing times per token, and \(\frac{l_j' + 1}{2}\) approximates the average decode-stage memory.

% \textbf{Equilibrium Conditions}: The batch processing rate matches arrivals:
% \[
% \frac{n_{j, \text{prefill}}^*}{\tau^*} = \lambda_j \quad \text{and} \quad \frac{n_{j, \text{decode}}^*}{\tau^*} = \lambda_j l_j' \quad \forall j,
% \]
% reflecting one prefill and \(l_j'\) decode iterations per prompt. Memory must satisfy:
% \[
% \sum_{j=1}^m \left( n_{j, \text{prefill}}^* l_j + n_{j, \text{decode}}^* \left( l_j + \frac{l_j' + 1}{2} \right) \right) \leq C,
% \]
% where \(C\) is the GPU memory capacity.

% \textbf{Throughput}: The equilibrium throughput remains:
% \[
% \throughput^* = \sum_{j=1}^m \lambda_j l_j',
% \]
% consistent with token generation rates.

% This refinement aligns the model with practical constraints, assuming a steady batch composition in equilibrium.

% \subsection{Equilibrium Throughput as a Benchmark}


% \arc{Add details on : 1. overloaded system and underloaded system 2. Lemma on upper bound of fluid throughput over all online non-preemptive policy $\Pi$. 3. Check the calculation.}
